% 図：コンパイラ型言語とインタプリタ型言語の実行の流れ（chapters/python/what_is_python.tex）
\begin{tikzpicture}
  % コンパイラ型（C 言語）
  \node[font=\small\bfseries, anchor=west] at (-1.1,1.25) {コンパイラ型（C 言語）};
  \node[figbox, minimum width=18mm] (csrc) at (0,0) {ソースコード\\{\footnotesize\ttfamily hello.c}};
  \node[figpart, minimum width=18mm] (cc) at (2.9,0) {コンパイラ\\{\footnotesize\ttfamily gcc}};
  \node[figbox, minimum width=18mm, fill=black!8] (bin) at (5.8,0) {実行ファイル\\{\footnotesize（機械語）}};
  \node[figbox, minimum width=14mm, fill=orange!10] (cpu1) at (8.5,0) {CPU};
  \draw[figarrow] (csrc) -- (cc);
  \draw[figarrow] (cc) -- node[figlabel, above]{翻訳} (bin);
  \draw[figarrow] (bin) -- node[figlabel, above]{実行} (cpu1);
  \node[figlabel, text=black!60] at (2.9,-0.95) {実行する前にまとめて翻訳};
  % インタプリタ型（Python）
  \node[font=\small\bfseries, anchor=west] at (-1.1,-1.85) {インタプリタ型（Python）};
  \node[figbox, minimum width=18mm] (psrc) at (0,-3.1) {ソースコード\\{\footnotesize\ttfamily hello.py}};
  \node[figpart, minimum width=47mm] (py) at (4.35,-3.1) {インタプリタ\\{\footnotesize\ttfamily python3}};
  \node[figbox, minimum width=14mm, fill=orange!10] (cpu2) at (8.5,-3.1) {CPU};
  \draw[figarrow] (psrc) -- (py);
  \draw[figarrow] (py) -- node[figlabel, above]{実行} (cpu2);
  \node[figlabel, text=black!60] at (4.35,-4.05) {1 行ずつ読みながら、その場で解釈して実行};
\end{tikzpicture}
